% Generated draft: requires mathematical and source review.
\documentclass[11pt]{article}
\usepackage[margin=25mm]{geometry}
\usepackage{fontspec}
\setmainfont{Latin Modern Roman}
\usepackage{amsmath,amssymb,mathtools,unicode-math}
\setmathfont{Latin Modern Math}
\usepackage{enumitem,needspace,xurl,hyperref}
\setlength{\parindent}{0pt}
\setlength{\parskip}{5pt}
\setlength{\emergencystretch}{3em}
\newcommand{\N}{\mathbb{N}}
\newcommand{\Z}{\mathbb{Z}}
\newcommand{\Q}{\mathbb{Q}}
\newcommand{\R}{\mathbb{R}}
\newcommand{\C}{\mathbb{C}}
\newcommand{\F}{\mathbb{F}}
\newcommand{\A}{\mathbb{A}}
\newcommand{\PP}{\mathbb{P}}
\newcommand{\E}{\mathbb{E}}
\newcommand{\eps}{\varepsilon}
\newcommand{\dd}{\,\mathrm{d}}
\begin{document}
\section*{UnsolvedMath ID: 20000310}
\textbf{Source problem code:} catalogue entry
\subsection{Definitions and mathematical statement}
Alexeev and Musta\c t\u a ask for a nontrivial application of motivic integration or arc spaces to termination of flips, and specifically ask whether such an application exists when every model in a flip sequence is smooth. A termination application should show finiteness of a stated class of sequences $X_1\dashrightarrow X_2\dashrightarrow\cdots$ using an invariant, bound or criterion obtained from arcs or motivic integration. The original does not specify the field, boundary, relative base or flip hypotheses. The isolated final ``1'' in the extraction has no stated mathematical role.

The attached unreviewed aid claims that Shokurov's termination criterion combined with the Ein--Musta\c t\u a--Yasuda arc-space lower-semicontinuity theorem proves termination of an ordinary boundaryless fixed-dimensional flip sequence whose every model is smooth, in the characteristic-zero setting to which those cited results apply. It describes the relevant minimal log discrepancies as integral arc-space codimensions in a finite set, making their ACC automatic. A minimal log discrepancy is an infimum of log discrepancies $A_X(E)=1+\operatorname{ord}_E(K_Y-f^*K_X)$ over the specified divisorial centers; the center convention and full termination-criterion hypotheses are not supplied.

For the narrower standard smooth projective flip of type $(m,\ell)$ with $m>\ell$, over a center $S$, the aid reports a gap $m-\ell$ in the codimensions of the maximal divisorial arc sets associated with the common exceptional valuation. ``Maximal divisorial'' refers to the arc set attached to that valuation, not a maximum over all valuations, whose discrepancies need not be bounded. It compares this with a published Chow-motive splitting removing $m-\ell$ Tate-twisted copies of the center motive. In a complex smooth projective interpretation let $B(V)=\sum_j\dim_\Q H^j(V(\C),\Q)$. It gives
\[
B(X)-B(X')=(m-\ell)B(S).
\]
For a chain of $N$ such standard flips the reported identities telescope to
\[
\sum_{i=1}^N(m_i-\ell_i)B(S_i)=B(X_1)-B(X_{N+1}),
\qquad N\le B(X_1)-1.
\]
The exact standard-flip geometry and published splitting are not reproduced. This quantitative Betti statement is limited to that standard smooth projective class and is not asserted for arbitrary flips or positive characteristic.
\subsection{Short English statement}
Use arcs or motivic integration to prove a genuine flip-termination statement. The aid proposes a smooth boundaryless application and a separate Betti-number budget for standard smooth projective flips, with their scopes kept distinct.
\subsection{Sources}
\begin{itemize}
\item Alexeev and Musta\c t\u a, Problem 7 in AIM \emph{Numerical invariants of singularities and higher-dimensional algebraic varieties}: \url{https://aimath.org/WWN/singularvariety/problems_final.pdf}. Canonical location: \texttt{aim-algebraic-geometry-notes.json}, notes[309]. Original wording is labeled exact despite extraction damage and a trailing isolated ``1''; remarks and literature are empty. No PDF was consulted.
\item The attached aid invokes Shokurov's criterion, Ein--Musta\c t\u a--Yasuda lower semicontinuity and an unnamed published Chow-motive splitting, without full citations. Its candidate quantitative corollary has low novelty confidence; the aid is machine generated and not peer reviewed.
\item Local registry: \texttt{data/unsolvedmath/problems.json}, permanent catalogue entry, exact source code \texttt{catalogue entry}.
\end{itemize}
\end{document}
